\section{Extending detected jets in selected coordinate directions}

\subsection*{Finding M-017: a combined second-row and later-row jet criterion}
Let \(r\geq2\), \(\ell\geq0\), \(t\geq0\), and
\(2\leq k\leq r+1\) be integers. Define
\[
 b=2\ell+2,\quad R=r-1,\quad p=k-2,\quad
 T=\lceil\sqrt R\rceil\binom{R+\ell}{\ell},\qquad
 N=T\left(1+\binom{p+t}{t}\right).
                                                               \tag{M-017.1}
\]
For every \(n\geq\max\{N,b+t\}\), the determinant
map on \(n\)-by-\(n\) linear pencils dominates the
coefficient projection with nonfirst-variable support
\[
 \Gamma=\left\{\gamma\in\mathbb Z_{\geq0}^{r-1}:
       |\gamma|\leq b+t,
       \ \sum_{i=k}^r\gamma_i\leq b\right\}.
                                                               \tag{M-017.2}
\]
The empty sum at \(k=r+1\) is zero. Consequently,
for every outer degree \(d\),
\begin{equation}\label{eq:M017K}
 \lambda\vdash nd,\quad\ell(\lambda)\leq r,\quad
 \lambda_2\leq b+t,\quad\lambda_k\leq b
       \quad\Longrightarrow\quad K_n(\lambda)=a_\lambda(d[n]).
\end{equation}
Missing parts of a partition are zero. When \(p\geq1\),
the constructed graph itself has a path of length \(b+t-1\),
so \(N\geq b+t\) and the extra maximum in the size
hypothesis is redundant. When \(p=0\), the target is just
the jet of order \(b\), independent of \(t\); in that
case \(n\geq N\) suffices.

\paragraph{An extension at each inward vertex.}
Start with the acyclic graph of M-013 at jet order \(b\).
It has \(T\) inward vertices and \(T\) outward vertices,
and its selected inward-to-outward inverse entries contain
every monomial of degree one through \(b-1\), each as
a single path product. All variables here are
\(x_1,x_2,\ldots,x_r\).

For each original inward vertex \(v\), make a separate
monomial tree on multi-indices \(\alpha\) in the first
\(p\) nonfirst variables \(x_2,\ldots,x_{k-1}\), with
\(|\alpha|\leq t\). Identify this new tree's root
with \(v\), and identify no other vertices. Orient the
new edges toward the root, label them by the variable
removed by the parent rule, and use positive labels.
There are \(\binom{p+t}{t}-1\) new vertices at each
old inward vertex. The resulting total number is exactly
\(N\) in (M-017.1).

The graph remains acyclic. A path beginning in a new tree
travels inward to its old vertex, then follows old inward
edges, crosses at most one old bridge, and travels outward.
No edge leads from the old graph into a new nonroot vertex.
In particular, if an old selected inverse entry is the
monomial \(x^\eta\), the entry from its new vertex
\(\alpha\) to the same outward endpoint is exactly
\[
                              x^\alpha x^\eta.
                                                               \tag{M-017.3}
\]
Uniqueness of the two path segments proves this identity;
the other attached trees introduce no alternate directed path.
For its adjacency matrix \(B\), use \(A_0=x_1I-B\).
Then \(\det A_0=x_1^N\), and after \(x_1=1\)
the selected determinant cofactors are the monomials in
(M-017.3), just as in M-013.

\begin{proof}[Coverage and dominant coefficient projection]
Every \(\gamma\in\Gamma\) with \(|\gamma|\leq b\)
is already a detected direction of M-013. Suppose instead
that \(b<|\gamma|\leq b+t\), and put
\(j=|\gamma|-b\). The suffix bound in (M-017.2)
implies that the total exponent in the early variables is
at least \(j\). Choose a degree-\(j\) submonomial
\(x^\alpha\) there. It has \(j\leq t\).
The remaining monomial has degree exactly \(b\).
Remove one variable from it for the linear entry variation
\(H\), leaving a monomial of degree \(b-1\geq1\).

The old M-013 graph represents that degree-\(b-1\)
monomial as an inward-to-outward inverse entry. Extend its
source by \(\alpha\) in the corresponding new tree.
Equation (M-017.3), multiplied by \(H(1,x)\), gives
exactly \(x^\gamma\) as the derivative of one matrix-entry
coefficient. Thus every target coefficient monomial is an
elementary derivative with coefficient one. Distinct target
monomials use distinct source coordinates, so the selected
Jacobian minor is the identity over \(\mathbb Z\).

The relevant paths also check the degree bound stated above.
If \(p\geq1\), take an early pure power of degree \(t\)
and an old pure-power path of degree \(b-1\).
Their concatenation is a simple path of length \(b+t-1\),
and uses \(b+t\) distinct vertices. Hence \(N\geq b+t\).
For \(p=0\), only the old jet target of degree \(b\)
occurs, and \(N=2T\geq b\).

Append an \(x_1I\) block when \(n>N\). Restricting
to the selected coefficient directions and translating the
target by its value at \(A_0\) yields \(F(t)=t+O(t^2)\)
over \(\mathbb Z\). The lowest homogeneous term proves
injectivity of polynomial pullback after any commutative
base change, and therefore dominance over \(\mathbb C\).

By F-016, a highest-weight polynomial of type \(\lambda\)
under the hypotheses of \eqref{eq:M017K} uses only
coefficient generators with total nonfirst degree at most
\(\lambda_2\leq b+t\) and suffix degree at most
\(\lambda_k\leq b\). It is a polynomial in the target
coordinates (M-017.2). If it lies in the determinant ideal,
it vanishes after the proved dominant projection and must
be zero. This argument applies to every highest-weight
copy, proving the complete ambient multiplicity equality.
\end{proof}

\paragraph{A further family of BIP necessary conditions.}
Define the explicit integer
\[
                  C_*=
       \left\lfloor\frac{190^{24}}{\binom{190^2+11}{12}}\right\rfloor.
                                                               \tag{M-017.4}
\]
Exact integer division gives \(C_*=478126755\).
At \(n\geq m^{25}\), every actual multiplicity obstruction
with \(\lambda_2\leq26+t\) satisfies
\begin{equation}\label{eq:M017BIP}
 \lambda_k\geq27
 \quad\text{whenever}\quad
 2\leq k\leq r+1,\qquad
                      1+\binom{k-2+t}{t}\leq C_*.
\end{equation}
This uses the known permanent-size reduction \(m\geq190\)
for an actual obstruction. It adds a coupled restriction
to the earlier single-row staircase; no source threshold
has been changed silently.

\begin{proof}
Positive permanent multiplicity gives \(r\leq m^2\).
Use (M-017.1) with \(b=26\), \(\ell=12\).
Since \(\lceil\sqrt{r-1}\rceil\leq m\),
\[
            T\leq m\binom{m^2+11}{12}.
\]
The product formula shows that
\(\binom{m^2+11}{12}/m^{24}\) is nonincreasing
for \(m\geq2\): up to the constant \(1/12!\)
it is \(\prod_{i=0}^{11}(1+i/m^2)\).
By definition of \(C_*\), for every \(m\geq190\)
one therefore has
\[
                  C_*\binom{m^2+11}{12}\leq m^{24}.
\]
If the binomial condition in \eqref{eq:M017BIP} holds,
the constructed size is at most \(m^{25}\leq n\).
The path argument checks any needed target degree bound.
If also \(\lambda_k\leq26\), the full multiplicity
equality just proved contradicts an actual obstruction.
This proves \eqref{eq:M017BIP}.
\end{proof}

For example, take \(k=3\). Then the binomial condition
is \(t+2\leq C_*\), so every BIP actual obstruction
must satisfy
\[
                     \lambda_3\geq27
                         \quad\text{or}\quad\lambda_2\geq C_*+25.
                                                               \tag{M-017.5}
\]
At the smallest allowed second row, \(\lambda_2=27\),
put \(t=1\). The condition becomes \(k\leq C_*\).
Thus such a candidate must have at least \(C_*\) rows,
and its rows \(2,\ldots,C_*\) all equal twenty-seven.
If it had fewer rows, taking \(k=r+1\) would contradict
\eqref{eq:M017BIP}. These are restrictions on remaining
candidates, not assertions that any such candidate is an
obstruction or even has positive permanent multiplicity.

\paragraph{Geometric dictionary and distinction from exact containment.}
The selected target is a finite torus-weight coefficient
module defined by two nested degree bounds. The prefix
extensions are integral affine-matrix morphisms, and the
unit Jacobian gives an actual universally injective coordinate
pullback. The relevant gradings are SGA~3, I,
Proposition~4.7.3, pp.~35--36; all Schur copies are
controlled by F-016's use of SGA~6, VI, Theorem~3.3,
pp.~392--393, and Macdonald, I(4.3), p.~63, Appendix~A
(5.3)--(5.6), pp.~155--156.
Higher coefficient coordinates remain free. Dominance of
the projection does not supply exact representations of
every form supported on \(\Gamma\), and does not follow
from a virtual \(K_0\) class. The strict positive-copy
example O-022 separately verifies that this distinction
can exclude BIP shapes missed by every exact support cover
of F-018 and every entire one-cut family.
